\subsection*{同一行处理中可能隐藏哪些事实}

重新取平台余额 10000、尚无收款。A 已提交 P1、C1、T1，W1 领取 T1。历史 \(H_0\) 中，W1 在发送以前退出。历史 \(H_1\) 中，平台已经提交 X1，余额成为 7000，W1 在保存 R 以前退出。两条历史的本地购买和支付记录相同，T1 也都留着 W1 的领取信息。

租期到期能够让 W2 继续工作，却没有给 \(H_0\) 和 \(H_1\) 增加支付结果证据。若到期就把 P1 改为付款失败，\(H_1\) 将出现平台已扣款、本地宣告拒绝的矛盾。若到期就按新身份收款，\(H_1\) 可以发生第二笔 3000 分扣款。接管的判断需要观察任务，发布付款结论则需要观察平台证据。

因此把本地状态分成三个维度。P1 的 \code{business\_state} 表示购买业务阶段。T1 的领取字段表示哪次工作者运行有资格修改调度。C1 的 \code{evidence\_state} 表示本地已核对到哪类平台结论。这些维度分别由受理与结论事务、领取与调度事务、可信证据合流事务更新。

\begin{table}[htbp]
\centering\small
\begin{tabularx}{\linewidth}{@{}p{.16\linewidth}p{.37\linewidth}X@{}}
\toprule
观察对象 & 典型取值 & 该事实回答的问题\\
\midrule
P1 业务阶段 & \code{awaiting\_payment}、\code{paid}、\code{payment\_rejected} & 这项购买已经保存了什么业务结论\\
T1 执行归属 & \code{ready}，或 W2 持有第 8 代，或 \code{done} & 哪次领取有资格继续安排恢复\\
C1 本地证据 & \code{unconfirmed}、\code{succeeded}、\code{rejected\_final} & 已保存证据能否确定原支付的结局\\
\bottomrule
\end{tabularx}
\caption{业务结论、调度资格与本地证据分别记录}\label{tab:ch08-dimensions}
\end{table}

平台在自己的账本中维护实际付款状态。本地的 \code{unconfirmed} 只表示尚缺确定结论，平台可能尚未收到请求、正在处理，或者已经成功。任务的到期时间同样要读成资格判断的一个条件：数据库可以继续保存一条过期的 \code{leased} 行，下一次领取才把它推进。调度入口读取持有者、代次和数据库时刻，合起来判断这次请求是否有权修改任务。

\CEightFigure{ch08-original02}{工作归属的变化与支付证据的变化各有依据}{fig:ch08-three-dimensions}

\subsection*{把允许的修改交给相应入口}

W1 取得一轮资格以后，可以续期、保存本轮诊断、安排下一次运行，或按适配规则准备查询与发送。这些修改会影响下一位工作者何时取得任务，所以调度入口检查它是否为当前持有者、是否携带当前代次、租期是否有效。平台观察则交给统一证据入口：该入口验证响应来源、原支付关联和固定输入，再决定能否进入 R。前一组检查授权谁安排工作，后一组检查授权哪份支付事实改变购买结论。

设 T1 原由 W1 持有，后来 W2 已接管。W1 恢复并请求把下一次运行时间改到明天。这个请求携带旧领取代次，调度入口拒绝它，否则 W1 会延后 W2 正在恢复的工作。W1 再带回一份平台成功响应，包含正确商户范围、C1、X1、CNY 3000 和原输入。证据入口验证其来源与关联后，可以把它交给同一个 R。

服务会重新校验原始响应的可信来源，或者通过平台查询复核，并逐项核对 C1 的固定输入、平台范围及编号。自报成功却缺少可验证来源的字符串，会在入口检查中被拒绝。证据记录保存来源、原始语义和处理决定。旧工作者的代次可以作为诊断资料留存，支付结论则由已核验的平台事实确定。

\CEightFigure{ch08-authority}{同一旧工作者的调度请求与支付证据，在各自入口分别核验}{fig:ch08-authority}

本章采用证据接收服务直接调用合流事务的路径。在接收请求返回成功以前，观察与相应 R 共同提交。进程在提交前退出，提交方可以重传，同一支付事实在 C1 处合流。若改为先持久收件再异步处理，收件及其待处理工作要共同保存，并由扫描器持续发现。这里的同步接收路径让经过授权与核验的回调、查询和旧工作者证据直接进入 R。

\subsection*{回调先完成时当前工作者怎样退出}

让平台从 10000 扣去 3000，W2 正在等待主动查询响应，T1 的第 8 代租期尚未到。可信回调先到，R 锁住 C1 与 T1，保存成功、O1 和原结果，把 T1 改为 \code{done}，清除持有者，并推进一次 \code{lease\_epoch}。这次推进撤销正在使用的领取版本。

随后 W2 试图续期或安排重试。它重新锁定并读取记录，看到 T1 已完成、代次已变，调度更新影响零行。W2 结束本轮。它较晚收到同一 X1 的成功，可以交给证据合流入口，入口恢复已保存 O1。收到较早的处理中，则记录旧观察并保留成功。

\CEightFigure{ch08-original03}{回调可以先提交 R，当前持有者随后读取完成事实}{fig:ch08-callback}

若 W2 的短调度事务先取得行锁，它可以先保存一次未来安排。R 随后取得锁并共同发布终态与 \code{done}，那次安排随之失去执行资格。若 R 先提交，W2 的修改被终态和代次检查拒绝。两个顺序都保留唯一 O1 和稳定支付结果。代码若只按 \code{task\_id} 更新、忽略已有 \code{done}，就可能把刚完成的 T1 再写回 \code{ready}。

R 提交结束的是这项本地恢复任务。W2 可能在它提交时等待网络响应，平台也可能继续处理此前已经到达的调用。查询接口按 P1 定位原 C1，读取已保存结果，调度器从 T1 读到 \code{done} 后停止派发。迟到的付款请求到达平台时，用 C1 已固定的身份恢复同一项支付。

\subsection*{先看 R 新增的任务完成步骤}

第 7.4 节已经给出支付证据的完整分支。本章增加两个受保护位置：入口锁序扩展为 C→T；首次保存支付成功或身份最终拒绝时，在提交 R 以前把同一 T 结束并撤销领取。处理中及仅未来关闭的观察保留任务，已有终态则核对任务已经完成。

\CEightFigure{ch08-r-delta}{R 的入口增加任务锁，首次终态把任务完成与支付结论共同提交}{fig:ch08-r-delta}

\subsection*{加入任务后的完整合流过程}

R 接收经过来源验证及字段核对的观察。事务先锁定实际 C，再按固定关联锁定 T，重核两行与 P 的归属，随后采样 \code{clock\_timestamp()}。缺失关联或发现未决 C 对应 \code{done}，整笔事务回滚并核查。第 7.4 节对重复、较早观察、内容冲突和身份最终拒绝的判断继续适用。任务结束动作与成功或拒绝写入共享同一事务：

\begin{BookCode}
已锁定实际 C、关联 T，读取 C 指向的 P，锁后采样数据库时刻
若 C 已终态：
    核对 T 已 done；核对重复证据或记录冲突；恢复原完整业务响应
若观察未决、只拒绝一次尝试，或仅关闭未来而既往未明：
    保存观察，保留非 done 的 T，提交后结束本次证据处理
若首次成功：
    按第 7.4 节插入订单，取 RETURNING 的真实 O
    构造完整 payment_result {purchase_id:P, charge_id:C, order_id:O,
        status:succeeded, provider_scope:原范围,
        provider_charge_id:已核对X, currency:原币种, amount_minor:原金额}
    条件更新 C，保存 succeeded、X、原平台结果及完整 payment_result
    条件更新 P，仅将 business_state 改为 paid；两次各要求一行
若首次身份最终拒绝：
    要求证据确认固定身份此前无成功，且尚无订单
    要求证据排除所有已受理、在途及未来尝试成功
    构造完整 payment_result {purchase_id:P, charge_id:C, status:rejected_final,
        reason:已核对原因, currency:原币种, amount_minor:原金额}
    条件更新 C，保存 rejected_final、原平台结果及完整 payment_result
    条件更新 P，仅将 business_state 改为 payment_rejected；两次各要求一行
上述任一首次终态分支继续：
    条件更新同一 T，要求 charge_id=C 且 phase 为 ready 或 leased
    令 phase=done、owner_incarnation=null、lease_until=null
    lease_epoch 加一，要求恰好返回一行 task_id、charge_id、lease_epoch
    保存观察，与订单、证据、结果和购买状态共同提交 R
\end{BookCode}

任务结束由支付证据授权。可信回调可以在 T 为 \code{ready} 时完成，也可以撤销 W2 的有效领取。返回零行表明关联或前提不符合当前迁移，整笔 R 回滚，订单、结果、观察与任务变化共同撤销，再按原 C 重读或保存冲突核查。主键、平台 X 唯一冲突以及其他约束异常采用第七章的回滚出口，代次溢出同样回滚。终态确认丢失后重读原 C/T：已提交则恢复原结果、保留已提交代次，尚未提交则按原证据执行。

完成后的响应使用第 7.4 节那三个完整对象：\code{acceptance} 固定，\code{progress} 来自当前购买 P，\code{payment\_result} 来自 C 保存的结果，三者在同一读取快照中取得。领取代次保存在内部任务记录中，用于恢复资格检查。对外结果版本以及完整购买身份和平台键，都取自原 P/C 的绑定。
